% Source: https://artofproblemsolving.com/community/c3h2524144_sigmas_problem_thread
% Stable identity: aops-topic-2524144-q50
From each of the $n$ vertices, diagonals go to the $n-3$ vertices other than itself and its two neighbours, and each diagonal is counted from both ends. So an $n$-gon has $\dfrac{n(n-3)}2$ diagonals. Then $n(n-3)=40$, which gives $n=8$ (the other root, $-5$, is not allowed). A $10$-gon has $\dfrac{10\cdot7}2=\boxed{35}$ diagonals.

\begin{asy}
usepackage("amsmath");
size(3.5cm);
defaultpen(fontsize(10pt));
pair[] V;for(int i=0;i<8;++i)V.push(dir(90+45*i));
for(int i=0;i<8;++i)for(int j=i+2;j<8;++j){if(i==0&&j==7)continue;draw(V[i]--V[j],gray(0.55));}
for(int i=0;i<8;++i)draw(V[i]--V[(i+1)%8],linewidth(0.9));
for(int i=0;i<8;++i)dot(V[i]);
\end{asy}

The octagon with its $20$ diagonals.
